% Generated draft: requires mathematical and source review.
\documentclass[11pt]{article}
\usepackage[margin=25mm]{geometry}
\usepackage{fontspec}
\setmainfont{Latin Modern Roman}
\usepackage{amsmath,amssymb,mathtools,unicode-math}
\setmathfont{Latin Modern Math}
\usepackage{enumitem,needspace,xurl,hyperref}
\setlength{\parindent}{0pt}
\setlength{\parskip}{5pt}
\setlength{\emergencystretch}{3em}
\newcommand{\N}{\mathbb{N}}
\newcommand{\Z}{\mathbb{Z}}
\newcommand{\Q}{\mathbb{Q}}
\newcommand{\R}{\mathbb{R}}
\newcommand{\C}{\mathbb{C}}
\newcommand{\F}{\mathbb{F}}
\newcommand{\A}{\mathbb{A}}
\newcommand{\PP}{\mathbb{P}}
\newcommand{\E}{\mathbb{E}}
\newcommand{\eps}{\varepsilon}
\newcommand{\dd}{\,\mathrm{d}}
\begin{document}
\section*{UnsolvedMath ID: 30003154}
\textbf{Source problem code:} catalogue entry
\subsection{Definitions and mathematical statement}
Let \(p\in\N\) tend to infinity through positive integers. The source denotes higher-codimension \(p\)-widths of \(M\) by \(\omega_p^k(M)\) and its Riemannian volume by \(\operatorname{vol}(M)\). The displayed exponents suggest that \(M\) has ambient dimension \(n+1\) and that the objects measured by \(\omega_p^k\) have dimension \(k\). These dimensional conventions are not explicitly defined in the supplied excerpt. Neither the admissible values of \(n,k\) nor the precise class of manifolds, boundary conditions, coefficient groups, cycle spaces, or sweepout classes defining these widths are specified. Consequently, the following records the mathematical targets with those qualifications unresolved.

\begin{enumerate}
\item Determine the constants denoted by \(a(n)\) in the source. The underlying Weyl-law normalization defining these constants is not supplied, and their relationship to the constants \(a(n,k)\) below is not stated explicitly.
\item Decide whether there are constants \(a(n,k)\), depending only on the indicated dimensions, such that, for every compact Riemannian manifold \(M\) in the intended but unspecified class and every admissible pair \((n,k)\) in the intended higher-codimension range,
\[
\lim_{p\to\infty}
\omega_p^k(M)\,p^{-\frac{n+1-k}{n+1}}
=
a(n,k)\,\operatorname{vol}(M)^{\frac{k}{n+1}}.
\]
The formula is preserved exactly as supplied. The compact Riemannian setting is identified in the supplied literature assessment; the excerpt does not specify whether manifolds are closed or may have boundary, or how boundary conditions enter the widths.
\item Determine analogous asymptotic limits for the quantities called \emph{Steenrod-power widths} in the source. The relevant cohomology operations, indexing, width construction, and asymptotic normalization are not supplied, so no specific limit formula for this part can be extracted.
\item Determine whether \(p\)-parameter min--max families yield minimal surfaces that become equidistributed as \(p\to\infty\). The source does not specify the dimension meant by ``surfaces,'' the quantifier over families and choices of min--max outputs, multiplicities, averaging procedure, or the measures and convergence defining equidistribution. This part therefore remains a qualitative question rather than a fully quantified convergence assertion.
\end{enumerate}

These are statement-only targets. The supplied open-status labels and literature assessment are dated source curation, not mathematical conclusions verified in this draft.
\subsection{Short English statement}
The source asks for Weyl-law constants, a higher-codimension width asymptotic on compact Riemannian manifolds, analogous Steenrod-power asymptotics, and equidistribution of min--max minimal surfaces. It supplies the displayed higher-codimension limit but omits several definitions, hypotheses, and convergence conventions needed for a complete precise formulation.
\subsection{Sources}
\begin{itemize}
\item Local registry: \texttt{data/unsolvedmath/problems.json}, supplied record for permanent UnsolvedMath ID \(30003154\), source problem code \texttt{catalogue entry}, titled ``Higher-Codimension Weyl Laws and Equidistributed Min--Max Surfaces.'' The record belongs to the Oberwolfach Reports Open Problems collection, gives proposed year 2016, and supplies no named proposer.
\item Supplied report citation: \emph{Geometrie (2017)}, Oberwolfach Reports, pp.~1577--1634, DOI \texttt{10.4171/owr/2016/28}; \url{https://doi.org/10.4171/owr/2016/28}. The background describes Oberwolfach Reports as records of mathematical workshops and their open-problem discussions. The supplied citation's 2017 label and the record's 2016 year information are retained without reconciliation.
\item Supplied literature reference: Liokumovich--Marques--Neves, \emph{Weyl law for volume spectrum}; \url{https://doi.org/10.4007/annals.2018.187.3.7}.
\item Supplied literature reference: \emph{Oberwolfach Report 28}; \url{https://doi.org/10.4171/owr/2016/28}.
\item Literature assessment dated \texttt{2026-08-21}, copied as supplied dated curation and not independently verified: status ``open.'' The assessment states that the Weyl law is proved for hypersurface volume spectra and for higher-codimension widths in Euclidean domains, while the corresponding compact Riemannian manifold problem and the stated Steenrod and equidistribution conclusions remain open. The background expressly describes this as dated literature triage derived from the supplied OWR curation and requests independent verification of time-sensitive status and cited claims.
\item The registry labels the statement \texttt{corrected\_verified}, explaining that notation, typography, and sentence structure were repaired without changing the mathematical content. This supplied editorial label does not establish mathematical correctness or resolve the missing formulation details. No external source was visited for this draft.
\end{itemize}
\end{document}
